\section{What is programmable cryptography?}
Cryptography is everywhere now and needs no introduction.
\emph{Programmable cryptography} is a term coined by 0xPARC for a second
generation of cryptographic primitives that has arisen in the last 15 or
so years.

To be concrete, let us consider two examples of what protocols designed
by classical cryptography can do:

\begin{description}
\item[Digital signatures.] RSA and ElGamal are examples of digital
  signature algorithms, where Alice can perform some protocol to prove
  to Bob that she endorses a message. A more complicated example might
  be a group signature scheme,\footurl{https://en.wikipedia.org/wiki/Group\_signature}
  which allows one member of a group to sign a message on behalf of the
  group.

\item[Confidential computing.] For example, consider
  Yao's millionaire problem,\footurl{https://en.wikipedia.org/wiki/Yao\%27s\_Millionaires\%27\_problem}
  where Alice and Bob want to know which of them makes more money
  without learning anything more about each other's incomes. With
  cryptography, Alice and Bob could use a two-party computation protocol
  designed specifically for this purpose.
\end{description}

Classically, first-generation cryptography relied on coming up with a
protocol for solving given problems or computing certain functions. The
goal of the second-generation ``programmable cryptography'' can then be
described as:

\begin{quote}
We want cryptography that can be ``programmed'' to work on
\textbf{arbitrary} problems and functions, rather than designing
protocols on a per-problem or per-function basis.
\end{quote}

To draw an analogy, it is like going from single-purpose hardware (like
a digital alarm clock or thermostat) to a general-purpose device (like a
smartphone) which can do any computation so long as someone writes code
for it.

\begin{remark}{}{}
The quote on the title page
(``I have a message $M$ such that $\sha(M) = \mathtt{0x91af3ac}\dots$'')
is a concrete example.
The hash function $\sha$ is a particular set of arbitrary instructions,
yet programmable cryptography promises that such a proof can be made
using a general compiler rather than inventing an algorithm specific to SHA-256.
\end{remark}

\section{Ideas in programmable cryptography}

Our work presents programmable cryptography through specific topics in
several self-contained ``easy pieces,'' imitating Richard Feynman's
wonderful approach to physics exposition. We quickly preview them here.

\subsection{2PC: two-party computation}

In a \alert{two-party computation (2PC)}, two people want to jointly
compute some known function \[F ( x_{1},x_{2}),\] where the
$i$-th person only knows the input $x_{i}$, without either person
learning the other person's input.

For example, in Yao's millionaire problem, Alice and Bob want to know
who has a higher income without revealing their own amounts. This is the
case where $F$ is the comparison function
($F ( x_{1},x_{2} )$ is $1$ if $x_{1} > x_{2}$, $2$ if
$x_{2} > x_{1}$, and $0$ if the two inputs are equal), and $x_{i}$
is the $i$-th person's income.

Two-party computation makes a promise that we will be able to do this
for \emph{any} function $F$ as long as we can implement it in code. It
generalizes to \alert{multi-party computation (MPC)}, which is one of the
main classes of programmable cryptography.

\subsection{SNARKs: proofs of general statements}

A powerful way of thinking about a signature scheme is that it is a
\textbf{proof}. Specifically, Alice's signature is a proof that ``I
[the person who generated the signature] know Alice's private key.''
Similarly, a group signature can be thought of as a succinct proof that
``I know one of Alice, Bob, or Charlie's private keys.''

In the spirit of programmable cryptography, a \emph{SNARK} generalizes
this concept as a ``proof system'' protocol that produces efficient
proofs of \textbf{arbitrary} statements of the form:

``I know $X$ such that $F(X,Y) = Z$, where $F,Y,Z$ are public,''
once the statement is encoded as a system of equations. One such
statement would be ``I know $M$ such that
$\sha(M) = Y$.''

SNARKs are an active area of research, and many different SNARKs are
known. We will focus on a particular example, PLONK (Section \ref{plonk}).

\subsection{FHE: fully homomorphic encryption}
Imagine you have some private text that you want to translate into
another language. While many services today will do this, even for free,
we can also imagine that you care about security a lot and you really do
not want the translating service to know anything about your text at all.

In \alert{fully homomorphic encryption (FHE)}, one person encrypts some
data $x$, and then a second person can perform arbitrary operations on
the encrypted data $x$ without being able to read $x$.

With this technology, you have a solution to your problem! You simply
encrypt your text $\Enc(x)$ and send it to your FHE
machine translation server. The server will faithfully translate it into
another language and give you $\Enc(y)$, where $y$ is
the translation of $x$. You can then decrypt and obtain $y$, knowing
that the server cannot extract anything meaningful from
$\Enc(x)$ without your secret key.

\subsection{ORAM: Oblivious RAM}
You want to perform a private computation on a large database. The
database is so large that you cannot store it yourself -- and you do not
trust the server it is stored on.

First off, you will encrypt the data, so the server cannot read it. But
the server still has an attack: They can study your \emph{access
patterns}. For example, they can see which records you access most
frequently, or which records you access at the same time as other
records. In many applications this is enough for the server to learn
sensitive information.

\alert{Oblivious RAM (ORAM)} protects against exactly this sort of attack.
Oblivious RAM is an algorithm you use to ``scramble'' your
memory access requests. When you feed your request into the ORAM
algorithm, the ORAM algorithm sends some scrambled read and write
requests to the server. Only one of the scrambled requests is the
request you are interested in; the others keep the server from learning
which request you care about.

\section{Programmable cryptography in the world}

In the past decade, there has been a surprisingly high amount of
theoretical work but also a surprisingly low amount of implementation
work on primitives in programmable cryptography. However, recent
advances in areas like blockchain and other decentralized systems are
rapidly driving demand for practical implementations of programmable
cryptography. The gap that is being revealed right now, as theory meets
reality, is exciting and enlightening.

Many of the protocols we mention in this book can be implemented today,
but only at a very high cost (for example, the cost of proving a
computation in a SNARK can be millions of times the cost of performing
the computation directly). As we study the theory of programmable
cryptography, it is useful to keep in mind some practical questions. Can
we reduce the theoretical overhead of programmable cryptography? How can
we make programmable cryptography systems more performant for modern
hardware and software systems? What other systems or applications can be
built on top of this technology?

It is easy to be carried away by the staggering possibilities, and to
imagine a perfect ``post-cryptographic'' world where everyone has
control over all their data and everyone's security preferences are
completely fulfilled. It is also easy to be cynical and assume that
these ideas will go nowhere. Reality is always somewhere in the middle;
the Internet today offers free search and civilization-scale
repositories of information to everyone, but is also used for plenty of
frivolous or even antisocial activity.

No matter what the future actually holds, one thing is clear - it is up
to people who are capable, curious, and optimistic to guide the next
stage of the evolution of cryptography-based systems. We hope that these
``easy pieces'' will inspire you to read, imagine, and build.
